% ============================================================
% CHƯƠNG 2: TÍNH KHẢ THI VÀ ĐỊNH RÕ VẤN ĐỀ
% ============================================================
\section{Tính Khả Thi Và Định Rõ Vấn Đề}

% ---- Slide 2.1: Phát biểu bài toán ----
\begin{frame}{Định Rõ Vấn Đề Thiết Kế}
  \begin{block}{Phát biểu bài toán}
    Thiết kế và chế tạo robot nhện 6 chân (hexapod) có khả năng di chuyển ổn định,
    thực hiện các kiểu di chuyển cơ bản, điều khiển từ xa qua WiFi,
    với chi phí phù hợp môi trường học thuật.
  \end{block}
  \vspace{0.3cm}
  \textbf{Yêu cầu chức năng cốt lõi:}
  \begin{itemize}
    \item Di chuyển tiến, lùi, xoay trái/phải; tránh vật cản
    \item Mỗi chân có 3 DOF: khớp Coxa, Femur, Tibia
    \item 18 servo được điều khiển đồng bộ theo thuật toán dáng đi
    \item Nhận lệnh từ xa qua WiFi (WebSocket)
  \end{itemize}
\end{frame}

% ---- Slide 2.2: Tính khả thi kỹ thuật ----
\begin{frame}{Tính Khả Thi Kỹ Thuật}
  \begin{columns}[T]
    \column{0.48\textwidth}
      \begin{exampleblock}{Cơ cấu cơ khí}
        Hexapod 3 khớp/chân -- phổ biến trong nghiên cứu\\
        Chi tiết in 3D FDM PLA/PETG -- đủ bền cho robot $< 2$ kg
      \end{exampleblock}
      \vspace{0.2cm}
      \begin{exampleblock}{Dẫn động}
        Servo MG996R -- moment 9--13 kgf.cm\\
        Đáp ứng tốt với tải trọng robot cỡ nhỏ
      \end{exampleblock}
    \column{0.48\textwidth}
      \begin{exampleblock}{Điều khiển}
        ESP32 + PCA9685 -- 32 kênh PWM qua I$^2$C\\
        IK hình học -- không yêu cầu tài nguyên lớn
      \end{exampleblock}
      \vspace{0.2cm}
      \begin{exampleblock}{Giao tiếp không dây}
        ESP32 WiFi SoftAP + WebSocket\\
        Web App (React) -- không cần cài thêm phần mềm
      \end{exampleblock}
  \end{columns}
\end{frame}

% ---- Slide 2.3: Tính khả thi kinh tế ----
\begin{frame}{Tính Khả Thi Kinh Tế}
  \centering
  \begin{tabular}{lccr}
    \toprule
    \textbf{Hạng mục} & \textbf{ĐV} & \textbf{SL} & \textbf{Ước tính (VNĐ)} \\
    \midrule
    Servo MG996R         & Cái   & 18  & 1.440.000 \\
    Vi điều khiển ESP32  & Cái   & 1   &   150.000 \\
    Module PWM PCA9685   & Cái   & 2   &   160.000 \\
    Vật liệu in 3D (PLA) & Cuộn  & 1--2 &   400.000 \\
    Pin LiPo 3S 11.1V    & Cái   & 1   &   300.000 \\
    Pin điều khiển 7.4V  & Cái   & 1   &   150.000 \\
    Cảm biến siêu âm     & Cái   & 1   &    30.000 \\
    \midrule
    \textbf{Tổng cộng}   &       &     & \textbf{\textasciitilde{}2.530.000} \\
    \bottomrule
  \end{tabular}
  \vspace{0.4cm}

  \alert{Phù hợp với ngân sách đồ án môn học cấp trường}
\end{frame}

% ---- Slide 2.4: Kế hoạch triển khai ----
\begin{frame}{Kế Hoạch Triển Khai (15 Tuần)}
  \centering
  \begin{tabular}{cl}
    \toprule
    \textbf{Tuần} & \textbf{Nội dung công việc} \\
    \midrule
    1--3   & Khảo sát tài liệu, xác định yêu cầu, chọn phương án thiết kế \\
    4--6   & Thiết kế cơ khí, xuất bản vẽ, tiến hành in 3D \\
    7--8   & Lắp ráp cơ khí, thiết kế và thi công mạch điện tử \\
    9--11  & Lập trình thuật toán IK và gait controller \\
    12--13 & Tích hợp hệ thống, thử nghiệm và hiệu chỉnh \\
    14--15 & Hoàn thiện sản phẩm, viết báo cáo và chuẩn bị bảo vệ \\
    \bottomrule
  \end{tabular}
\end{frame}

% ---- Slide 2.5: Ràng buộc và tiêu chí ----
\begin{frame}{Ràng Buộc Và Tiêu Chí Thiết Kế}
  \begin{columns}[T]
    \column{0.48\textwidth}
      \textbf{Ràng buộc kích thước:}
      \begin{itemize}
        \item Thân (không tính chân): $\leq 25$ cm
        \item Rải rộng khi chân mở: $\leq 50$ cm
        \item Tổng trọng lượng: $< 2$ kg
      \end{itemize}
      \vspace{0.2cm}
      \textbf{Ràng buộc dẫn động:}
      \begin{itemize}
        \item Servo RC: moment $\geq 9$ kgf.cm
        \item Góc khai thác: $\pm 60°$ quanh trung tính
      \end{itemize}
    \column{0.48\textwidth}
      \textbf{Tiêu chí hiệu năng:}
      \begin{itemize}
        \item Tốc độ tối thiểu: 5 cm/s
        \item Hoạt động liên tục $\geq 15$ phút
        \item Tần số cập nhật servo $\geq 50$ Hz
      \end{itemize}
      \vspace{0.2cm}
      \textbf{Tiêu chí an toàn:}
      \begin{itemize}
        \item Cầu chì bảo vệ nguồn
        \item Servo về vị trí an toàn khi mất tín hiệu
      \end{itemize}
  \end{columns}
\end{frame}
